\subsection{Adding Physics: Physics-Informed PCE}

\begin{frame}[more-space]{Model\enspace|\enspace Adding Physics: Physics-Informed PCE}
  \dataheading{A simple thermal-resistance model as physical guidance}
  \begin{beamercolorbox}[sep=2mm,wd=\textwidth]{data panel}
    \footnotesize
    At equilibrium, temperature rise is approximated by heat generation
    multiplied by an effective thermal resistance:
    \[
      T_{\mathrm{phys}}=T_{\mathrm{inlet}}+\Delta T_0
        +P_{\mathrm{CPU}}R_{\mathrm{th}}(F),
      \qquad F=P_{\mathrm{fan,raw}}>0,
    \]
    \[
      R_{\mathrm{th}}(F)=R_0+K\left(\frac{F}{F_{\mathrm{ref}}}\right)^{-\gamma}.
    \]
  \end{beamercolorbox}
  \medskip
  \begin{columns}[T,onlytextwidth]
    \column{0.48\textwidth}
      \dataheading{Meaning of the parameters}
      \footnotesize
      $R_0\geq0$: resistance remaining at strong cooling.\par
      $K\geq0$: additional resistance at reference cooling.\par
      $\gamma>0$: how resistance decreases with cooling.\par
      $\Delta T_0$: fitted temperature offset.\par
      \smallskip
      $R_0,K$ have units of ${}^\circ$C/W.
    \column{0.49\textwidth}
      \dataheading{Calibration}
      \footnotesize
      Fix $F_{\mathrm{ref}}>0$ on the same raw scale so $F/F_{\mathrm{ref}}$
      is dimensionless. It is not a conversion to watts.\par
      \smallskip
      Calibrate the remaining parameters using training cases only.
      With 19 observations, check identifiability and simplify if needed.
  \end{columns}
\end{frame}

\begin{frame}[more-space]{Model\enspace|\enspace Adding Physics: Physics-Informed PCE}
  \begin{beamercolorbox}[sep=2mm,wd=\textwidth]{data panel}
    \footnotesize
    The prediction keeps the same orthonormal Legendre expansion.
    Only the objective used to learn its coefficients changes:
    \[
      \mathcal L(\mathbf c)=
      \underbrace{\frac{1}{N}\sum_{i=1}^{N}
        \left[\widehat T_{\max}(\mathbf x_i)-y_i\right]^2}_{\mathcal L_{\mathrm{data}}}
      +\lambda_{\mathrm{phys}}
      \underbrace{\frac{1}{M}\sum_{j=1}^{M}
        \left[\widehat T_{\max}(\mathbf z_j)-T_{\mathrm{phys}}(\mathbf z_j)\right]^2}
        _{\mathcal L_{\mathrm{physics}}}.
    \]
  \end{beamercolorbox}
  \medskip
  \begin{columns}[T,onlytextwidth]
    \column{0.48\textwidth}
      \dataheading{Fit measurements and encourage physics}
      \footnotesize
      $N$: training observations; $y_i$: measured targets.\par
      $M$: physics-check operating points $\mathbf z_j$.\par
      \smallskip
      Restrict checks to supported, validated operating conditions; generated
      points are not additional experimental evidence.
    \column{0.49\textwidth}
      \dataheading{Choose the physics influence}
      \footnotesize
      $\lambda_{\mathrm{phys}}\geq0$ controls the compromise.
      Zero gives the data-only baseline. Select it within training-case validation.\par
      \smallskip
      This is a loss-based idea similar to a PINN, using polynomials instead of a neural network.
  \end{columns}
\end{frame}

\begin{frame}[more-space]{Model\enspace|\enspace Adding Physics: Physics-Informed PCE}
  \begin{columns}[T,onlytextwidth]
    \column{0.48\textwidth}
      \begin{beamercolorbox}[sep=2mm,wd=\linewidth]{data panel}
        \dataheading{What it could improve}
        \footnotesize
        \begin{itemize}
          \item Encourage realistic heating and cooling trends.
          \item Potentially improve predictions with sparse data.
          \item Keep fast PCE, UQ and Sobol calculations.
        \end{itemize}
      \end{beamercolorbox}
    \column{0.49\textwidth}
      \begin{beamercolorbox}[sep=2mm,wd=\linewidth]{data warning}
        \dataheading{What it cannot guarantee}
        \footnotesize
        \begin{itemize}
          \item A penalty does not enforce exact physics.
          \item Missing input combinations remain untested.
        \end{itemize}
      \end{beamercolorbox}
  \end{columns}
  \medskip
  \begin{beamercolorbox}[sep=2mm,wd=\textwidth]{data panel}
    \dataheading{How we would decide whether it helps}
    \footnotesize
    \begin{itemize}
      \item Compare PCE and PI-PCE on the same held-out cases.
      \item Fit physics parameters and tune the loss weight using training cases only.
      \item Keep the simpler model unless accuracy and physical behavior improve.
      \item Extra tuning requires more validation.
    \end{itemize}
  \end{beamercolorbox}
\end{frame}
